\section{Recursive weaving and the universal bound}\label{sec:recursion}

The recursive use of woven districts follows the framework of
Postle \cite[Section~6]{Postle2020} and Delcourt and Postle
\cite[Section~7]{DelcourtPostle2021}. We now use separability to replace
the successive small steps of the preceding section by a recursion with
a fixed shrinkage factor.

\begin{theorem}\label{thm:recursive-woven}
There are absolute positive integers $K,Q$ such that, for every integer
$a\ge1$, every $Ka$-connected graph $P$ with $\ell(P)\ge Qa$ is
$(a,3a)$-woven.
\end{theorem}
\begin{proof}
The main inductive step concerns $h(P)<64a$ and sufficiently large $a$.
It separates the paths for prescribed pairs from the clique-model
construction. Give each root and each end of a noncoincident pair its
own neighboring proxy outside the original terminal set, with all proxies
distinct, and delete the original terminals. Choose a minimum-total-length
system of two paths per proxy to a small highly connected district $D$.
Deleting the proxies, $D$, and these paths costs only $O(a)$ in list
chromatic number. We can therefore extract a highly connected region
$H'$ disjoint from the first path system. Further connections
from $H'$ to $D$ and Lemma~\ref{lem:combine} produce disjoint paths from
all proxies and from a set $W$ of $a$ vertices of $H'$ to $D$, meeting
$H'$ only at $W$. It then suffices to build a clique model rooted at $W$
inside $H'$: grouping the path ends in $D$ attaches the model to the
root proxies and completes the prescribed pair paths.

For this internal model, divide the labels of $W$ into three blocks.
Separability and connectivity extraction supply three disjoint
districts for recursive calls, one for each pair of blocks. A call
needs at most $2\lceil a/3\rceil$ roots. Each label receives two bags
to be joined to its vertex of $W$. The $2a$ joining paths may pass through
the recursive districts; wovenness lets us reroute them while constructing
the child models, keeping the retained bags disjoint from the paths
except at their designated roots. Every pair of labels shares a district
and hence receives an adjacency.

We specify the order of the constants first. Fix all constants in the
preceding sections, including $L,C_*,B_*,K_0,D_{K_0}$, the deletion
constants $c_d$, and the separability constant $S$. Choose an integer
$K\ge100$ such that
\begin{equation}\label{recursive:K}
 K-14>128D_{K_0}.
\end{equation}
Choose an absolute integer $a_0\ge24$ so large that all sufficiently-large
parameter conclusions used below hold with $t=64a$ whenever $a>a_0$.
Enlarge $a_0$ if necessary so that
\begin{equation}\label{recursive:shrink}
 a'=2\lceil a/3\rceil\le3a/4\qquad(a>a_0),
\end{equation}
and $t^{1.06}+7a\le t^{11/10}$ for $t=64a$, $a>a_0$.
All these requirements involve only constants already fixed, and none
involves $Q$. Finally choose an integer $Q$ satisfying
\begin{align}
 Q/4&>10+64(c_{42}+6K_0+2S)+(9/2)K,
       \label{recursive:Q-induction}\\
 Q&>10+64(c_{42}+6K_0+3S),
       \label{recursive:Q-separation}\\
 Q&>128C_*\log(128a_0)+1.
       \label{recursive:Q-base}
\end{align}
We prove the assertion by strong induction on $a$.

\paragraph{Separating the terminal roles.}
Fix $a$ distinct roots and at most $3a$ pairs as in
Definition~\ref{def:woven}. Let $Z_0$ be their union, so $|Z_0|\le7a$.
For each root, and for each end of each noncoincident pair, choose a
neighbor outside $Z_0$ as its \emph{proxy}. All proxies are distinct,
even for different roles of the same original vertex. This is possible
greedily because $\delta(P)\ge Ka$ and fewer than $14a$ vertices are
forbidden at any choice. Coincident pairs will ultimately use trivial
paths at their original vertices and need no proxies.

Put $G=P-Z_0$ and let $Z$ be the proxy set. Then
\begin{equation}\label{recursive:proxy-bounds}
 a\le|Z|\le7a,\qquad
 G\text{ is }(K-7)a\text{-connected},\qquad
 \ell(G)\ge Qa-7a.
\end{equation}
It suffices to find in $G$ a clique model rooted at the root proxies and
paths for the other proxy pairs, with all model bags and all these paths
pairwise disjoint. Adding the proxy edges and the original vertices then
restores precisely the root--path coincidences allowed by wovenness.
Different prescribed pairs have no common original endpoint; a
coincident pair at a root meets only that root's expanded bag.

\paragraph{A large minor and the finite base range.}
Set $t=64a$. If $h(P)\ge t$, delete from a $K_t$ model every bag meeting
$Z_0$. This leaves a $K_{57a}$ model in $G$. Since
$57a\ge4|Z|$ and $(K-7)a\ge2|Z|$, Lemma~\ref{lem:rooting} gives a
clique model rooted at all of $Z$. Retain the root-proxy bags. For each
other proxy pair, a path through its two bags and an edge between them
joins the pair. These paths and the retained bags are disjoint because
all proxies and their bags are distinct. This proves the required
conclusion in the large-minor case, for every $a$.

Suppose $h(P)<64a$. Lemma~\ref{lem:coarse-density} gives
\[
 \ell(P)\le128C_*a\log(128a)+1.
\]
For $1\le a\le a_0$ this is strictly less than $Qa$ by
\eqref{recursive:Q-base}. Thus the minor-free case is impossible
throughout the base range. We may henceforth assume $a>a_0$.

\paragraph{A junction district and a fresh region.}
Every vertex of $G-Z$ has degree at least $(K-14)a$, which is greater
than $2D_{K_0}t$ by \eqref{recursive:K}. In particular this graph is
nonempty and has density greater than $D_{K_0}t$. Apply
Lemma~\ref{lem:district} to find a $K_0t$-connected district
$D\subseteq G-Z$ with $|V(D)|\le t^{1.06}$.

The connectivity of $G$ and the order of $D$ permit two paths per proxy
from $Z$ to distinct vertices of $D$, with only the prescribed repeated
starts in common, by Lemma~\ref{lem:double-packing}. Choose such a system
$\mathcal Q_1$ of minimum total length in $G$. It has at most $14a$ paths.
Put
\[
 X=Z\cup V(D),\qquad Y=V(\mathcal Q_1)\setminus X.
\]
We have $|X|\le t^{11/10}$ by the choice of $a_0$.
Lemma~\ref{lem:path-minimization} gives $28a$-degeneracy of $G[Y]$ and
at most $42a$ neighbors in $Y$ for every vertex outside $X\cup Y$.
These bounds are at most $42t$, so Lemma~\ref{lem:additive-deletion}
with parameter $42$ gives
\begin{equation}\label{recursive:after-paths}
 \ell(G-X-Y)\ge Qa-7a-c_{42}t.
\end{equation}
By \eqref{recursive:Q-separation} the right-hand side exceeds $6K_0t$.
Lemma~\ref{lem:list-extraction} therefore supplies a $K_0t$-connected
induced subgraph $H'\subseteq G-X-Y$ such that
\begin{equation}\label{recursive:Hprime}
 \ell(H')\ge Qa-7a-(c_{42}+6K_0)t.
\end{equation}

In $G-Z$, the sets $H'$ and $D$ have order greater than $K_0t$, and the
connectivity is at least $(K-14)a>2a$. Hence there are $2a$ disjoint
$H'$-to-$D$ paths. Truncate them at their last visit to $H'$ and their
first subsequent visit to $D$, so they have no internal vertex in either
set. They avoid $Z$. Lemma~\ref{lem:combine}, applied to these paths and
$\mathcal Q_1$, gives disjoint paths to $D$ from all proxies in $Z$ and
from a set $W$ of $a$ distinct vertices of $H'$. These paths meet $H'$ and
$D$ only at their appropriate ends.

It is now enough to construct a clique model rooted at $W$ entirely
inside $H'$. Indeed, in $D$ group the terminal ends of the outside paths
as follows: pair each root proxy with a different member of $W$, and
pair the other proxies as prescribed. There are at most $|Z|+a\le8a$
terminals, so Theorem~\ref{thm:grouping} applies because
$K_0t\ge8\gamma a$. These disjoint groups and the outside paths attach
the rooted model to the root proxies and supply all other required
paths. The model inside $H'$ has no unwanted intersection with the
outside paths, since their only vertices in $H'$ are precisely $W$.

\paragraph{Three recursive districts.}
For each $w\in W$ choose two neighbors in $H'-W$, all $2a$ choices
distinct. Minimum degree at least $K_0t$ permits this greedy choice.
Let $N$ be their set and put $H''=H'-W-N$. Then
\begin{equation}\label{recursive:r0}
 \ell(H'')\ge r_0:=Qa-10a-(c_{42}+6K_0)t.
\end{equation}
Equation~\eqref{recursive:Q-separation} gives $r_0>3St$.
Apply Lemma~\ref{lem:separability} to $H''$, and then to one of the two
resulting subgraphs. The first application is legal because
$\ell(H'')\ge r_0>2St$. Each of its subgraphs has list chromatic number
at least $r_0-St>2St$, making the second application legal as well.
We obtain three disjoint subgraphs each of list chromatic number at
least $r_0-2St$.

Take $a'=2\lceil a/3\rceil$. After the loss of $2St$ to separability,
each subgraph must still pay $6Ka'$ for connectivity extraction and
retain $Qa'$ for the recursive call. If
$A=10+64(c_{42}+6K_0+2S)$, then
\begin{align*}
 r_0-2St-6Ka'-Qa'
 &=Q(a-a')-Aa-6Ka'\\
 &\ge\bigl[Q/4-A-(9/2)K\bigr]a>0
\end{align*}
by \eqref{recursive:shrink} and \eqref{recursive:Q-induction}.
In particular $r_0-2St>6Ka'$, and list extraction in each subgraph
produces disjoint $Ka'$-connected induced districts $J_1,J_2,J_3$
with list chromatic number at least $Qa'$. Since $a'<a$, induction
makes each district $(a',3a')$-woven.

Partition the labels of $W$ into three nonempty blocks $B_1,B_2,B_3$,
each of size at most $\lceil a/3\rceil$. Assign one district to each
pair of blocks, relabeling them $J_{12},J_{13},J_{23}$ accordingly.
In its district, designate a distinct root for each
label of the corresponding two blocks, and add arbitrary distinct
padding roots to bring the total to $a'$. The districts have enough
vertices for this choice. Every label has exactly two designated roots,
in different districts; every pair of labels occurs together in at
least one district, including two labels in the same block; see
Figure~\ref{fig:three-districts}.

\begin{figure}[tb]
\centering
\begin{tikzpicture}[font=\small,
 block/.style={draw,rounded corners,minimum width=2.15cm,minimum height=.75cm,fill=blue!4},
 district/.style={draw,rounded corners,minimum width=2.15cm,minimum height=.75cm,fill=green!5},
 every path/.style={line width=.6pt}]
\node[block] (b1) at (0,2) {$B_1$};
\node[block] (b2) at (3.8,2) {$B_2$};
\node[block] (b3) at (7.6,2) {$B_3$};
\node[district] (j12) at (0,0) {$J_{12}: B_1\cup B_2$};
\node[district] (j13) at (3.8,0) {$J_{13}: B_1\cup B_3$};
\node[district] (j23) at (7.6,0) {$J_{23}: B_2\cup B_3$};
\draw[-{Latex[length=2mm]}] (b1.south) -- (j12.north);
\draw[-{Latex[length=2mm]}] (b1.south east) -- (j13.north west);
\draw[-{Latex[length=2mm]}] (b2.south west) -- (j12.north east);
\draw[-{Latex[length=2mm]}] (b2.south east) -- (j23.north west);
\draw[-{Latex[length=2mm]}] (b3.south west) -- (j13.north east);
\draw[-{Latex[length=2mm]}] (b3.south) -- (j23.north);
\end{tikzpicture}

\caption{The three recursive districts. The blocks $B_1,B_2,B_3$ partition
the labels, and a label has one designated root in each of the two
districts receiving its block. Arrows indicate label assignments, not
graph edges. Every pair of labels is represented together in a district.}
\label{fig:three-districts}
\end{figure}

Inside $H'-W$, join the two chosen neighbors of each $w$ to its two
designated roots, respectively, using disjoint paths. This is an
application of Theorem~\ref{thm:grouping} to $4a$ distinct terminals:
the neighbors lie in $N$, while the districts were chosen in $H''$
and hence avoid $N$. The connectivity is at least $K_0t-a\ge4\gamma a$.
There are $2a$ paths. Apply Lemma~\ref{lem:rerouting} successively in the
three districts to incorporate their rooted clique models while
retaining these path ends. Each district has sufficient weaving
capacity, since
\[
 2a\le3a'=6\lceil a/3\rceil.
\]
Discard all padding bags. Each designated root remains an endpoint
of its path; the retained bags meet the rerouted paths only at their
own designated roots. Since the districts are disjoint, later rerouting
does not introduce an intersection in a previously retained bag.

For each $w$, unite its two designated bags with their respective paths,
the two edges to $w$, and $w$ itself. The resulting $a$ sets are connected
and disjoint, rooted at $W$. Two labels have an edge between their sets
because they occur together in one of the district clique models.
This constructs the required rooted $K_a$ model inside $H'$ and
completes the induction.
\end{proof}

\begin{proof}[Proof of Theorem~\ref{thm:main}]
Let $K,Q$ be as in Theorem~\ref{thm:recursive-woven}. We first show that,
for every integer $t\ge2$,
\begin{equation}\label{recursive:final-minor-free}
 h(G)<t\quad\Longrightarrow\quad\ell(G)\le(Q+6K)t.
\end{equation}
Otherwise Lemma~\ref{lem:list-extraction}, with connectivity parameter
$Kt$, gives a $Kt$-connected induced subgraph $P$ with
$\ell(P)\ge\ell(G)-6Kt>Qt$. Theorem~\ref{thm:recursive-woven} makes
$P$ $(t,3t)$-woven. Its order is greater than $Kt\ge t$, so we may choose
$t$ roots and no pairs. Wovenness then supplies a $K_t$ model,
contradicting $h(G)<t$.

For a nonempty graph, $h(G)\ge1$. Taking $t=h(G)+1\le2h(G)$ in
\eqref{recursive:final-minor-free} proves
\[
 \ell(G)\le2(Q+6K)h(G).
\]
Thus the absolute integer $C=2(Q+6K)$ has all the claimed uniformity.
The finite base range was covered before the induction began, so no
small-parameter exception remains.
\end{proof}
